\section{Related Work}
\label{sec:related-work}

\paragraph{Long-term memory benchmarks.}
LoCoMo and LongMemEval evaluate long conversations, temporal updates, and multi-session retrieval \citep{maharana2024locomo,wu2025longmemeval}. MEME adds questions about deletion, cascade effects, and absent facts \citep{jung2026meme}. MemoryAgentBench (MAB) combines accurate retrieval, temporal and test-time learning, long-range understanding, and selective forgetting in one evaluation framework \citep{hu2025memoryagentbench}. These benchmarks measure different parts of the storage and reading problem, so we use them as complementary tests rather than reducing memory to one score.

\paragraph{Memory systems for agents.}
MemGPT treats context management as paging between a working window and an external store \citep{packer2023memgpt}. Model-Native Computing Architecture likewise frames context capacity, cache reuse, and agent scheduling as system-level resource-management problems \citep{lin2026icam}. Mem0, Zep/Graphiti, HippoRAG-2, A-MEM, MIRIX, and Cognee use summaries, notes, graphs, or combinations of them to persist information across sessions \citep{chhikara2025mem0,rasmussen2025zep,gutierrez2025hipporag2,xu2025amem,wang2025mirix,markovic2025cognee}. These systems differ in what they keep as the memory object: Deep-Dream keeps the source document and treats concepts and relations as links back to it, so a derived graph entry is a search handle rather than a replacement fact.

\paragraph{Temporal and versioned knowledge.}
Temporal knowledge graphs represent changing entities and time-dependent relations, while probabilistic record linkage separates false merges from missed matches \citep{fellegi1969theory}. Deep-Dream applies these ideas to agent memory: processing time orders observations and versions, and delayed alignment allows a candidate identity to remain revisable while its source evidence stays intact. Event dates remain part of the observations and can guide retrieval without creating a second state clock.

\paragraph{Compression, retrieval, and evidence.}
Information bottleneck and rate--distortion theory describe the trade between a compact representation and information needed for a target \citep{tishby2000information,shannon1959coding}. Associative-memory models study retrieval from partial cues \citep{ramsauer2021hopfield}, while retrieval-augmented generation selects external evidence for generation \citep{lewis2020rag}. Long-context models can lose relevant information when it is buried in a large window \citep{liu2024lost}; attribution work checks whether a response cites a source \citep{gao2023alce}. Deep-Dream combines these ideas in a persistent setting: concepts and relations narrow the search, source spans retain the wording, and the agent chooses how much evidence to load; \cref{sec:information-view} formalizes this separation.

\paragraph{Agentic retrieval.}
A-RAG exposes keyword, semantic, and chunk-reading tools to a reasoning model for multi-hop retrieval \citep{du2026arag}. Deep-Dream shares the idea that a model can choose its next retrieval action, but applies it to a changing personal history. Its additional requirements are identity revision, version links, cross-document episodes, and source-grounded verification.
